% Generated draft: requires mathematical and source review.
\documentclass[11pt]{article}
\usepackage[margin=25mm]{geometry}
\usepackage{fontspec}
\setmainfont{Latin Modern Roman}
\usepackage{amsmath,amssymb,mathtools,unicode-math}
\setmathfont{Latin Modern Math}
\usepackage{enumitem,needspace,xurl,hyperref}
\setlength{\parindent}{0pt}
\setlength{\parskip}{5pt}
\setlength{\emergencystretch}{3em}
\newcommand{\N}{\mathbb{N}}
\newcommand{\Z}{\mathbb{Z}}
\newcommand{\Q}{\mathbb{Q}}
\newcommand{\R}{\mathbb{R}}
\newcommand{\C}{\mathbb{C}}
\newcommand{\F}{\mathbb{F}}
\newcommand{\A}{\mathbb{A}}
\newcommand{\PP}{\mathbb{P}}
\newcommand{\E}{\mathbb{E}}
\newcommand{\eps}{\varepsilon}
\newcommand{\dd}{\,\mathrm{d}}
\begin{document}
\section*{UnsolvedMath ID: 20003117}
\textbf{Source problem code:} catalogue entry
\subsection{Definitions and mathematical statement}
The extracted source asks only whether \(\delta(\Gamma_n)\to\delta(\Gamma_0)\), attributed in its typography to F. Gu\'ertitaud. It does not include the definitions of the groups, their convergence, or \(\delta\). The verification field says the preceding source paragraph is essential but does not reproduce it. The following is the specific metric-ball subgroup interpretation recovered by the supplied aid.

Let \(\Gamma_0\) be a torsion-free discrete subgroup of \(\operatorname{Isom}^+(\mathbb H^2)\), let \(S=\mathbb H^2/\Gamma_0\), and fix \(x\in S\). Using the covering identification \(\pi_1(S,x)\cong\Gamma_0\), define, for integers \(n\geq1\),
\[
 \Gamma_n=\operatorname{im}\bigl(\pi_1(B(x,n),x)\longrightarrow\pi_1(S,x)\bigr).
\]
These are image subgroups; inclusion of the ball need not be \(\pi_1\)-injective. If torsion is allowed, an orbifold formulation is needed instead. The aid's exhaustion has
\[
 \Gamma_1\leq\Gamma_2\leq\cdots\leq\Gamma_0,
 \qquad\bigcup_{n\geq1}\Gamma_n=\Gamma_0,
\]
since each based loop has compact image in some finite-radius ball. The ambient notation \(\Gamma_0\) is not an initial subgroup in the increasing stage sequence.

For a subgroup \(H\leq\Gamma_0\), fix \(o\in\mathbb H^2\) and define its Poincar\'e series and critical exponent by
\[
 P_H(s)=\sum_{h\in H}e^{-s\,d(o,ho)},\qquad
 \delta(H)=\inf\{s\geq0:P_H(s)<\infty\}.
\]
The aid invokes Sullivan's 1979 Corollary 6 for the increasing-union equality
\[
 \delta(\Gamma_0)=\sup_{n\geq1}\delta(\Gamma_n).
\]
Subgroup inclusion gives monotonicity of the stage exponents, so the reported conclusion is
\[
 \delta(\Gamma_n)\nearrow\delta(\Gamma_0).
\]
This conclusion uses the cited critical-exponent theorem, not just formal monotone convergence of positive series; the complete hypotheses and reference must be verified.

The supplementary formulation applies to any nested exhaustion by subsets whose carried image subgroups \(H_n\) contain every compact based loop at some stage, so \(\bigcup_nH_n=\Gamma_0\). It reports the same exponent limit. Moreover, at \(\delta=\delta(\Gamma_0)\), monotone convergence of the sums gives
\[
 P_{\Gamma_0}(\delta)=\lim_nP_{H_n}(\delta).
\]
Therefore the ambient group is of divergence type, meaning \(P_{\Gamma_0}(\delta)=\infty\), exactly when these stage sums are unbounded, allowing an individual stage sum already to be infinite. This does not assert convergence or divergence type from the exponents alone. The aid describes the recovered metric-ball question as solved affirmatively; without its missing source context, no general statement about an arbitrary sequence \(\Gamma_n\) is intended.
\subsection{Short English statement}
Under the aid's recovered ball-image interpretation, the groups form an increasing union, and the cited Sullivan theorem makes their critical exponents increase to the ambient exponent. A separate monotone-sum criterion detects divergence type at the limit. The definitions absent from the extracted source remain explicit reconstruction qualifications.
\subsection{Sources}
\begin{itemize}
\item Local UnsolvedMath registry, permanent ID \(20003117\), exact source code \texttt{catalogue entry}; canonical \texttt{aim-topology-notes.json}, \texttt{notes[204]}, source tag \texttt{problem}.
\item AIM workshop ``Lipschitz metric on Teichmueller space,'' Problem 4.4, attribution extracted as F. Gu\'ertitaud: \url{https://aimath.org/pastworkshops/lipschitzteichproblems.pdf}. The record labels the statement exact, but the supplied verification notes essential preceding context not included in the registry sentence. The PDF was not fetched, and the author's typography requires checking.
\item The aid cites Sullivan, 1979, Corollary 6, for the increasing-union critical-exponent theorem. No paper title, publication details, identifier, or full authoritative citation is supplied. The original literature field is empty.
\item Research attempt 1: result \texttt{full\_solution}, problem status at run \texttt{solved}, rigor label \texttt{proved}, novelty confidence \texttt{low}; classification \texttt{SOLVED-IN-LITERATURE}, registry status \texttt{solved}. Created and updated 14 August 2026. These are supplied assessments for the recovered interpretation, not independent verification.
\item Attribution: Ulam AI / UnsolvedMath Contributors, CC BY 4.0. The machine-generated research aid is not peer reviewed and requires independent verification before citation.
\end{itemize}
\end{document}
